\section{From an annular height to a flexible sphere}
\label{sec:geometric-closure}

This section proves that Theorem~\ref{thm:annulus-intro} implies
Theorem~\ref{thm:main}. It uses no analysis beyond the classical
theory of surfaces, and it fixes exactly what the remaining sections
have to construct. The steps are as follows.
\begin{itemize}
\item An isometric embedding with prescribed height $z$ exists locally
 if and only if the remainder metric $g-dz\otimes dz$ is positive and
 flat (Lemma~\ref{lem:darboux-invariant}).
\item On a disk whose rim is planar, a positive flat remainder metric
 develops onto a round disk, and the development is the identity near
 the rim (Proposition~\ref{lem:disk-development}).
\item The resulting graph is closed by a cap which carries a marker
 (Lemma~\ref{lem:fixed-cap}).
\item Turning the graph inside the cap gives the family of
 Theorem~\ref{thm:main} (Theorem~\ref{thm:geometric-closure}).
\end{itemize}

\subsection{Setting}

Throughout the section $0<r_-<r_+$, and on the annulus
$(r_-,r_+)\times\T$ with coordinates $(r,\theta)$ we consider a metric
\begin{equation}\label{eq:intro-metric}
 g=E(r)\,dr^2+r^2\,d\theta^2,\qquad E=1+w^2,
\end{equation}
where $w$ is a smooth function of $r$. This is the metric induced on
the graph $(r\cos\theta,r\sin\theta,h(r))$ of a radial function with
$h'=w$. We identify $(r,\theta)$ with the point
$x=(r\cos\theta,r\sin\theta)$ of the plane, and write $D_\rho$ for the
closed disk of radius $\rho$ centred at the origin.

\begin{definition}\label{def:flat-endpoint}
A smooth function $f$ on $(r_-,r_+)$, or on $(r_-,r_+)\times\T$, is
\emph{flat} at an endpoint if it extends smoothly by zero across that
endpoint; equivalently, every derivative of $f$ tends to zero there. A
function $z$ has a \emph{constant extension} at an endpoint if $z$ has
a limit $c$ there, independent of $\theta$, and $z-c$ is flat.
\end{definition}

If $w$ is flat and $z$ has constant extensions at both endpoints, we
may extend $w$ by zero and $z$ by its limits, and regard them as
defined on a slightly larger annulus. On the added collars $w=0$ and
$z$ is constant. All statements of this section are therefore
formulated for functions which are constant on open collars of the two
ends.

\subsection{Removing one coordinate from the metric}

A map $F=(P,z)$ into $\R^2\times\R$ is isometric for $g$ exactly when
\[
 P^*g_{\R^2}=\bar g,\qquad \bar g:=g-dz\otimes dz .
\]
Thus $\bar g$ must be a positive definite flat metric, and for an
embedding its planar development $P$ must be injective. Flatness is a
scalar equation for $z$. For a symmetric $2$-tensor $B$ we write
$\det_gB=\det(B_{ij})/\det(g_{ij})$, and $\nabla_g^2z$ is the Hessian
of $z$ for the metric $g$.

\begin{lemma}[Darboux equation]\label{lem:darboux-invariant}
Let $z$ be a smooth function on a surface with metric $g$, and assume
$\sigma:=|dz|_g^2<1$. Then $\bar g=g-dz\otimes dz$ is positive
definite, and its curvature is
\begin{equation}\label{eq:darboux-curvature}
 K_{\bar g}=
 \frac{K_g(1-\sigma)-\det_g\nabla_g^2z}{(1-\sigma)^2}.
\end{equation}
In particular $\bar g$ is flat if and only if
\begin{equation}\label{eq:intro-darboux}
 \det\nolimits_g\nabla_g^2z=K_g\,(1-|dz|_g^2).
\end{equation}
\end{lemma}

\begin{proof}
In a $g$-orthonormal coframe $(e^1,e^2)$ with $dz=\sqrt\sigma\,e^1$ we
have $\bar g=(1-\sigma)\,e^1\otimes e^1+e^2\otimes e^2$. This is
positive definite, and $|dz|_{\bar g}^2=\sigma/(1-\sigma)$.

Consider the product $(\Sigma,\bar g)\times\R$ and in it the graph
$x\mapsto(x,z(x))$. The metric induced on the graph is
$\bar g+dz\otimes dz=g$. Its upward unit normal has vertical component
$(1+|dz|_{\bar g}^2)^{-1/2}=\sqrt{1-\sigma}$. The vertical coordinate
is a function on the product with vanishing Hessian, and its restriction
to the graph is $z$; hence the intrinsic Hessian of $z$ is the second
fundamental form $\mathrm{II}$ of the graph times the vertical
component of the normal,
\[
 \nabla_g^2z=\sqrt{1-\sigma}\;\mathrm{II}.
\]
The sectional curvature of the product on the tangent plane of the
graph is $K_{\bar g}$ times the squared area distortion of the
horizontal projection, which is $1-\sigma$. The Gauss equation for the
graph therefore reads
\[
 K_g=(1-\sigma)K_{\bar g}+\det\nolimits_g\mathrm{II}
    =(1-\sigma)K_{\bar g}+(1-\sigma)^{-1}\det\nolimits_g\nabla_g^2z ,
\]
and solving for $K_{\bar g}$ gives \eqref{eq:darboux-curvature}.
\end{proof}

Equation~\eqref{eq:intro-darboux} is thus exactly the compatibility
condition for the two horizontal coordinates once the vertical
coordinate has been chosen. It guarantees local developments. The
global reconstruction, including injectivity, is the subject of the
next proposition; no local embedding theorem is used to bypass it.

\subsection{Development onto a round disk}

A flat metric on an annulus may have nontrivial holonomy around the
hole. Filling the inner collar by a flat disk removes the holonomy,
and the half-turn symmetry of $z$ forces the inner disk to stay
centred.

\begin{proposition}[Planar reconstruction]\label{lem:disk-development}
Suppose that $w$ vanishes on open collars of both ends of
$(r_-,r_+)$, that $z$ equals a constant $c_L$ on the inner collar and
a constant $c_R$ on the outer one, and that
$z(r,\theta+\pi)=z(r,\theta)$. Suppose that $\bar g=g-dz\otimes dz$ is
positive definite and flat. Extend $w$ by $0$ and $z$ by $c_L$ to
$r\le r_-$. Then there is an orientation-preserving diffeomorphism
$P:D_{r_+}\to D_{r_+}$ such that
\[
 P^*g_{\R^2}=\bar g,
 \qquad P=\id\ \text{on the outer collar},
 \qquad P=R_\gamma\ \text{on the inner disk }D_{r_-}
\]
for some angle $\gamma$. Consequently $F=(P,z):D_{r_+}\to\R^3$ is an
isometric embedding of $(D_{r_+},g)$, and its image is a graph over
$D_{r_+}$.
\end{proposition}

\begin{proof}
\step{1: a normalized development}
On the inner disk and on the outer collar we have $w=0$ and $dz=0$, so
there $\bar g=g=dr^2+r^2d\theta^2$ is the Euclidean metric in the
coordinate $x$. In particular $\bar g$ is a smooth flat metric on the
whole disk $D_{r_+}$, including its centre. Choose a global oriented
$\bar g$-orthonormal coframe $\eta_1,\eta_2$ which equals $dx_1,dx_2$
on the outer collar, and let $\alpha$ be its connection form:
$d\eta_1=\alpha\wedge\eta_2$, $d\eta_2=-\alpha\wedge\eta_1$. Flatness
means $d\alpha=0$, and since the disk is simply connected,
$\alpha=d\psi$ with $\psi=0$ on the outer collar. The complex form
$e^{i\psi}(\eta_1+i\eta_2)$ is then closed. Its primitive $P$,
normalized to be the identity on the outer collar, is an
orientation-preserving local isometry from $(D_{r_+},\bar g)$ to the
plane.

\step{2: the inner disk is rotated about its centre}
Two orientation-preserving local isometries from a connected flat
surface to the plane differ by a fixed Euclidean motion; this follows
by differentiating their relative rotation. Let $I(x)=-x$ be the
half-turn. It preserves $\bar g$ because $g$ is rotationally symmetric
and $z\circ I=z$. Both $P$ and $-P\circ I$ are orientation-preserving
local isometries equal to the identity on the outer collar, so they
coincide: $P$ is odd. On the inner disk $\bar g$ is Euclidean, hence
$P$ is a Euclidean motion there, and an odd Euclidean motion is a
rotation $R_\gamma$ about the origin.

\step{3: $P$ is a diffeomorphism of the disk}
Extend $P$ by the identity outside $D_{r_+}$. The extension is a local
diffeomorphism of $\R^2$ which is the identity outside a compact set;
it is proper, hence a covering map, and it has one sheet because it is
injective near infinity. So it is a diffeomorphism of $\R^2$, and
since it fixes the exterior of the disk it maps $D_{r_+}$ onto itself.

\step{4: conclusion}
The map $F=(P,z)$ is injective because $P$ is, and its induced metric
is $P^*g_{\R^2}+dz\otimes dz=\bar g+dz\otimes dz=g$. Its image is the
graph of $z\circ P^{-1}$ over $D_{r_+}$.
\end{proof}

\begin{figure}[htbp]
\centering
\begin{tikzpicture}[scale=.87,>=Stealth,font=\small]
\begin{scope}
\filldraw[fill=blue!6,draw=black] (0,0) circle (1.65);
\filldraw[fill=green!12,draw=green!45!black] (0,0) circle (.58);
\draw[dashed] (0,0) circle (1.42);
\draw[->] (0,0)--(1.25,.55) node[right] {$r$};
\node[below,align=center] at (0,-1.9) {intrinsic disk\\metric $g-dz\otimes dz$};
\end{scope}
\draw[->,thick] (2.1,0)--(3.9,0) node[midway,above] {$P$};
\node[align=center] at (3,-.65) {flat,\\no holonomy};
\begin{scope}[xshift=6cm]
\filldraw[fill=blue!6,draw=black] (0,0) circle (1.65);
\filldraw[fill=green!12,draw=green!45!black] (0,0) circle (.58);
\draw[dashed] (0,0) circle (1.42);
\draw[->] (.1,.1) arc[start angle=45,end angle=210,radius=.25];
\node at (0,.95) {$P=\id$ near the rim};
\node[below,align=center] at (0,-1.9) {planar image:\\the same round disk};
\end{scope}
\end{tikzpicture}
\caption{The development $P$. Filling the inner collar removes the
holonomy of the annulus; the normalization fixes the outer collar; the
half-turn symmetry keeps the inner disk centred, so that $P$ only
rotates it. Properness upgrades the local development to a
diffeomorphism. The sketch records these constraints, not the
pointwise distortion of $P$.}
\label{fig:development}
\end{figure}

\subsection{A stationary cap which distinguishes shapes}

The graph of Proposition~\ref{lem:disk-development} is horizontal near
its boundary, so it can be closed by a cap. A rotationally symmetric
cap would hide the relative rotation: the closed surfaces obtained for
different $\tau$ would be congruent. We therefore put on the cap a
marker which is recognized by every congruence.

\begin{lemma}[The marked cap]\label{lem:fixed-cap}
Let $S_0$ be a smooth graph over $D_{r_+}$ which is horizontal near its
boundary. Then $S_0$ extends to a smooth embedded sphere $S_0\cup S_1$.
The complementary cap $S_1$ can be chosen so that, for any given
constant $K_{\max}$, it contains a point $p$ with
\[
 K(p)>K_{\max},\qquad K(p)>K(p')\ \text{ for every other point }p'
 \text{ of }S_1,
\]
and with distinct principal curvatures at $p$.
\end{lemma}

We shall take $K_{\max}$ to be the maximum of the Gaussian curvature
of $S_0$. Then $p$ is the unique point of the closed surface at which
the Gaussian curvature is largest.

\begin{proof}
\step{1: the unmarked cap}
Translate so that the boundary height is zero, and extend the
horizontal collar to a radius $R>r_+$. Choose a smooth nondecreasing
function $\chi:[0,1]\to[0,1]$ which is $0$ near $0$, is $1$ near $1$,
and satisfies $\chi(1-t)=1-\chi(t)$. For $L>0$ consider the unit-speed
meridian curve $s\mapsto(r(s),z(s))$, $0\le s\le L$, which starts at
$(R,0)$ and satisfies
\[
 r'(s)=\cos(\pi\chi(s/L)),\qquad
 z'(s)=-\sin(\pi\chi(s/L)).
\]
Its tangent turns from the outward horizontal direction to the inward
one. By the symmetry of $\chi$ the curve stays at radius at least $R$
and returns to radius $R$, and its height decreases by a positive
constant times $L$. Take $L$ so large that the curve ends below the
minimum of the graph $S_0$, and complete it by the horizontal segment
from its endpoint to the axis. Rotating this meridian gives a smooth
embedded surface which meets $S_0$ only along the attaching collar.

\step{2: the marker}
Place a bump in the horizontal annulus $r_+<r<R$. In Cartesian
coordinates $(x,y)$ centred at a point of that annulus, choose
constants $A,B>0$ with $A\ne B$, and $k,\rho>0$; put $\xi=Ax^2+By^2$
and
\[
 u(x,y)=\frac{k}{2}\int_\xi^{\rho^2}
       e^{-s/(\rho^2-s)}\,ds\quad(\xi<\rho^2),
 \qquad u=0\quad(\xi\ge\rho^2).
\]
This function is smooth. With $\Xi=\xi/(\rho^2-\xi)$ one computes
$u_x=-kAx\,e^{-\Xi}$ and $u_y=-kBy\,e^{-\Xi}$, and the Gaussian
curvature of the graph of $u$ is
\begin{equation}\label{eq:marker-curvature}
 K=\frac{ABk^2e^{-2\Xi}\,(1-2\xi\,\Xi')}
 {\bigl[1+k^2e^{-2\Xi}(A^2x^2+B^2y^2)\bigr]^2},
 \qquad \Xi'=\frac{d\Xi}{d\xi}>0 .
\end{equation}
At the centre $K=ABk^2$. At every other point of the support the
numerator is either nonpositive or strictly smaller than $ABk^2$, and
the denominator is at least one. So the centre $p$ is the unique
maximum of $K$ on the bump, with value $ABk^2$. First choose $k$ so
that $ABk^2$ exceeds both $K_{\max}$ and the curvature of the unmarked
cap. Then choose $\rho$ so small that the support $\{\xi<\rho^2\}$
fits inside the horizontal annulus; the bump has height at most
$k\rho^2/2$ and does not meet the rest of the surface. At $p$ the
gradient of $u$ vanishes and its Hessian is $\diag(-kA,-kB)$, so the
principal curvatures are $-kA$ and $-kB$ for the upward normal. They
are distinct.
\end{proof}

\begin{figure}[htbp]
\centering
\begin{tikzpicture}[x=.8cm,y=.8cm,>=Stealth,font=\small]
\draw[thick,blue!70!black]
 (-3.3,0)--(-2.8,0) .. controls (-2.5,.1) and (-2.5,.9) .. (-1.5,.75)
 .. controls (-1,.65) and (-.8,.65) .. (-.5,.65)--(.5,.65)
 .. controls (.8,.65) and (1,.65) .. (1.5,.75)
 .. controls (2.5,.9) and (2.5,.1) .. (2.8,0)--(3.3,0);
\draw[thick,gray!75] (3.3,0)--(4,0)
 .. controls (5.5,0) and (5.5,-2) .. (4,-2)--(-4,-2)
 .. controls (-5.5,-2) and (-5.5,0) .. (-4,0)--(-3.3,0);
\fill[red!75!black] (3.7,0) circle (2.5pt);
\draw[dashed,gray] (-2.8,0)--(-2.8,-2);
\draw[dashed,gray] (2.8,0)--(2.8,-2);
\node[above] at (0,1.4) {constant inner disk};
\draw[->] (0,1.3)--(0,.7);
\node[above] at (-3.3,.3) {flat collar};
\node[above] at (3.8,.65) {marker};
\draw[->,red!75!black] (3.8,.6)--(3.7,.1);
\node[below] at (0,-2.1) {fixed cap, below the entire moving graph};
\end{tikzpicture}
\caption{A meridional sketch of the closed surface. The blue graph
depends on the angle; all its rotated copies stay inside the central
cylinder and above the bottom of the cap. The marker lies in the
stationary horizontal extension of the collar.}
\label{fig:geometric-closure}
\end{figure}

\subsection{The annulus implies the theorem}

\begin{theorem}[Geometric closure]\label{thm:geometric-closure}
Let $g$ be a metric of the form \eqref{eq:intro-metric} and let $z$ be
a smooth function on the annulus. Suppose that
\begin{enumerate}
\item $w=0$ on open collars of both ends, and $z$ is constant on each
 of these collars;
\item $z(r,\theta+\pi)=z(r,\theta)$, and $z$ is not independent of
 $\theta$;
\item $g-dz\otimes dz$ is positive definite and flat.
\end{enumerate}
Then all the conclusions of Theorem~\ref{thm:main} hold.
\end{theorem}

\begin{proof}
\step{1: the family}
Let $P$ and $F=(P,z)$ be given by
Proposition~\ref{lem:disk-development}, and attach to the graph
$F(D_{r_+})$ the marked cap of Lemma~\ref{lem:fixed-cap}, with
$K_{\max}$ the maximum of the Gaussian curvature of $F(D_{r_+})$. Let
$\Sigma$ be the resulting sphere, with its induced metric $g_\Sigma$.
On the disk put
\[
 F_\tau=R_{-2\pi\tau}\circ F\circ\varrho_{2\pi\tau},\qquad
 \varrho_\vartheta(r,\theta)=(r,\theta+\vartheta),
\]
and let $F_\tau$ be the inclusion on the cap. The intrinsic rotation
$\varrho_\vartheta$ preserves $g$, so each $F_\tau$ is isometric on
the disk. On the outer collar $F(x)=(x,c_R)$, and on the inner disk
$F(x)=(R_\gamma x,c_L)$; in both cases the two rotations cancel and
$F_\tau=F$. Hence $F_\tau$ agrees with the inclusion near the rim and
is smooth on $\Sigma$, jointly in the point and in $\tau$, and
$F_{\tau+1}=F_\tau$. The image $F_\tau(D_{r_+})=R_{-2\pi\tau}F(D_{r_+})$
is a graph over the same round disk with the same range of heights, so
it meets the cap only along the collar: $F_\tau$ is an embedding. The
inner disk and the cap are fixed pointwise.

\step{2: the disk images are not all equal}
Suppose $R_\vartheta F(D_{r_+})=F(D_{r_+})$ for every $\vartheta$. Then
$F^{-1}\circ R_\vartheta\circ F$ is an isometry of $(D_{r_+},g)$, and
on the outer collar it equals $\varrho_\vartheta$. An isometry of a
connected Riemannian manifold is determined by its value and
differential at one point, so $F^{-1}R_\vartheta F=\varrho_\vartheta$
on the whole disk. Comparing the vertical components of
$R_\vartheta F=F\varrho_\vartheta$ gives
$z(r,\theta+\vartheta)=z(r,\theta)$ for all $\vartheta$, contrary to
assumption~(ii).

\step{3: the closed images are not all congruent}
The Gaussian curvature of $F_\tau(\Sigma)$ is determined by
$g_\Sigma$, so by the choice of the marker the point $p$ is, for every
$\tau$, the unique point of $F_\tau(\Sigma)$ of largest Gaussian
curvature; it lies on the cap and is the same point of $\R^3$ for all
$\tau$. Let $\Phi$ be a Euclidean isometry mapping $F_0(\Sigma)$ onto
$F_\tau(\Sigma)$. It preserves the Gaussian curvature, so
$\Phi(p)=p$. It maps the bounded component of the complement to the
bounded component, so it preserves the outward normal at $p$, and then
it preserves each of the two principal directions at $p$, which are
well defined because the principal curvatures are distinct and which
are the same for all $\tau$ because the cap is fixed. There are at
most four such isometries. Thus, if all the images were congruent, the
family $\tau\mapsto F_\tau(\Sigma)$ would be a continuous family of
compact sets taking values in a finite set, hence constant. The part
of $F_\tau(\Sigma)$ inside the closed cylinder over $D_{r_+}$ and above
the bottom of the cap is $F_\tau(D_{r_+})$, so the disk images would
all be equal, which Step~2 excludes.

\step{4: the invariants}
The area is constant because the metric is. The moving part
$F_\tau(D_{r_+})$ is a rotated copy of $F(D_{r_+})$ and the rest of
the surface is stationary, so the integral of the oriented mean
curvature is constant. Finally let $Z_\tau$ be the height of
$F_\tau(D_{r_+})$ as a graph over $D_{r_+}$, and let $c_B$ be the
height of the bottom of the cap. The part of the enclosed volume
inside the cylinder over $D_{r_+}$ is
\[
 \int_{D_{r_+}}(Z_\tau(y)-c_B)\,dy
 =\int_{D_{r_+}}(Z_0(R_{2\pi\tau}y)-c_B)\,dy
 =\int_{D_{r_+}}(Z_0(y)-c_B)\,dy,
\]
and the remaining volume is bounded by the stationary part of the
surface.
\end{proof}

\begin{proof}[Proof that Theorem~\ref{thm:annulus-intro} implies
Theorem~\ref{thm:main}]
Let $w$ and $z$ be as in Theorem~\ref{thm:annulus-intro}. Extend $w$
by zero and $z$ by its two limiting constants to a slightly larger
annulus. The extended functions are smooth, and on the added collars
$g-dz\otimes dz=g$ is Euclidean, hence positive and flat. All the
hypotheses of Theorem~\ref{thm:geometric-closure} hold.
\end{proof}

What remains is therefore precise: construct one nonradial height on
an annulus with a positive flat remainder metric and with constant
extensions at both ends. The family of embeddings then follows without
any further existence argument.
